\section{Balancing parts and zeros of the Ehrhart numerator}
\label{sec:numerator}

\begin{proposition}[Schur concavity]\label{prop:schur}
At fixed vertex count $v$, moving a vertex from a larger part to a
smaller one to make the sizes more equal weakly increases every
$\gamma_j$. Equivalently, the gamma coefficients are Schur-concave in
the partition of $v$, with zero parts allowed for comparison.
\end{proposition}
\begin{proof}
Fix $j\ge1$ and put $d=v-1$. In \eqref{eq:multipartite-gamma} the loss
for one part of size $a$ is
$F(a)=\sum_{\ell>j}\binom a\ell\binom{d-a}{2j-\ell}$.
Adding one vertex to the marked part newly disqualifies precisely the
sets containing this vertex, $j$ old marked vertices and $j-1$ old
unmarked vertices. Thus
\[
 D(a):=F(a+1)-F(a)=\binom a j\binom{d-a-1}{j-1}
 \qquad(0\le a<d).
\]
Each factor is log-concave in $a$ with interval support, so their
product is unimodal. Consider the inverse move $(a,b)\mapsto(a+1,b-1)$
for $a\ge b\ge1$, with both graphs connected. Set $u=b-1$, so
$u\le a\le d-u$ and $a<d$. If $D(u)>0$, then $u>0$ and
$D(d-u)/D(u)=(d-u)/u\ge1$. Unimodality on $[u,d-u]$ gives
$D(a)\ge D(u)$. The inequality also holds when $D(u)=0$. Hence the
loss increases by $D(a)-D(u)\ge0$. Such unit transfers generate
majorization of integer partitions, and $\gamma_0=1$ throughout.
\end{proof}

Deleting $M$ disjoint edges from $K_{d+1}$ gives $M$ parts of size two,
the remaining parts being singletons. Formula~\eqref{eq:multipartite-gamma}
specializes to the gamma polynomial
\begin{equation}\label{eq:gamma-match}
 \Gamma_{d,M}(z)=G_d(z)-2Mz,\qquad
 G_d(z)=\sum_{j=0}^{\lfloor d/2\rfloor}\binom{2j}{j}\binom d{2j}z^j .
\end{equation}
The real-rootedness of $h^*$ for every complete bipartite symmetric edge
polytope~\cite[Thm.~4.8]{hjm} therefore raises a different question from
the location of the roots of $E_G$ itself.

\begin{theorem}[Threshold for real roots of the numerator]\label{thm:nonnumerator}
Write $N_d$ for the number of nonreal roots of $\Gamma_{d,M}$ with
multiplicity. The number for $h^*$ is exactly $2N_d$. Set
\[
 A_0(t)=\frac{e^{2t}}{2\sqrt\pi\,t^{5/4}},\qquad
 c_{\rm num} = A_0(5/8)=\frac{e^{5/4}}{2\sqrt\pi}(8/5)^{5/4}.
\]
\begin{enumerate}
\item For $d\ge4$ there is a real threshold $T_d$ with
$T_d/d^{3/4}\to c_{\rm num}$, such that $N_d=0$ if and only if $M\le T_d$.
\item If $M/d^{3/4}\to c>c_{\rm num}$, let $t_-<t_+$ solve $A_0(t)=c$.
Then $N_d/\sqrt d\to 2(\sqrt{t_+}-\sqrt{t_-})/\pi$.
At $c=c_{\rm num}$, $N_d=o(\sqrt d)$.
\item If $M/d^\alpha\to\lambda>0$ with $3/4<\alpha\le1$, subject to
$2M\le d+1$, then
\[
 \frac{N_d}{\sqrt{d\log d}}\ \longrightarrow\ 
 \frac{2}{\pi}\sqrt{\frac{\alpha-3/4}{2}} .
\]
In particular, a maximum matching gives
$(\sqrt2/\pi+o(1))\sqrt{d\log d}$ nonreal roots of $h^*$.
\end{enumerate}
\end{theorem}

\begin{proof}
Put $P(x)=G_d(-x)$ and $m=\lfloor d/2\rfloor$.
Let $L_d$ denote the Legendre polynomial normalized by $L_d(1)=1$.
Summing the binomial series in \eqref{eq:gamma-match} over $d$ gives
$\sum_d P(x)z^d=(1-2z+(1+4x)z^2)^{-1/2}$.
The Legendre generating function therefore gives
\[
 P(x)=(1+4x)^{d/2} L_d(\cos\theta),\qquad
 \theta=\arctan(2\sqrt{x}),\qquad x>0 .
\]
In particular $P$ has $m$ distinct positive roots, from the positive
zeros of $L_d$. We use the usual Legendre identities, orthogonality and
normalization~\cite[Ch.~18]{dlmf}.

\emph{Counting horizontal intersections.}
The negative roots of $\Gamma_{d,M}$ correspond to the positive
solutions of $Y(x)=M$, where $Y(x)=-P(x)/(2x)$. Its critical points
are given by $xP'(x)-P(x)=0$. This polynomial has at least one
zero in every gap between consecutive positive roots of $P$, by Rolle
applied to $Y$. It also has exactly one negative zero, since its value
at $x=-z$, $z>0$, is $-1+\sum_{j\ge2}(j-1)
\binom{2j}{j}\binom d{2j}z^j$. Its degree $m$ accounts for all
these zeros, each simple. Thus a bounded positive lobe of $Y$ between
consecutive zeros has one maximum. At $M>0$ it contributes two
solutions with multiplicity precisely when this maximum is at least $M$.
The first and last exterior intervals are monotone, and any positive
terminal branch attains every positive level once.
All gamma coefficients are nonnegative by Theorem~\ref{thm:gamma}, so
$\Gamma_{d,M}$ has no nonnegative real root. It follows that $N_d$ is
twice the number of positive bounded lobes whose maxima are below $M$.
For $m\ge2$ such positive lobes exist, and $T_d$ is the smallest of
their maxima. At $M=0$ all roots are real as well. Each nonreal root
$z$ of the gamma polynomial produces two distinct nonreal roots of
$z(1+t)^2=t$ in $h^*$. Each real gamma root is negative and gives
two real solutions. Any additional numerator zeros are at $t=-1$.

\emph{Legendre amplitude.}
Set $n=d+\tfrac12$, $u(\theta)=\sqrt{\sin\theta}\,L_d(\cos\theta)$,
and $q=n^2+1/(4\sin^2\theta)$. The Legendre equation gives
$u''+q u=0$. On $(0,\pi/2)$, $q$ decreases, so for
$S=u^2+u'^2/q$ we have $S'=-q'u'^2/q^2\ge0$ and $(qS)'=q'u^2\le0$.
Consequently
\[
 \frac{q(\pi/2)}{q(\theta)}S(\pi/2)\le S(\theta)\le S(\pi/2),
 \qquad S(\pi/2)\sim\frac{2}{\pi d}.
\]
The last estimate follows from
$L_{2h}(0)=(-1)^h\binom{2h}{h}/4^h$, $L_{2h+1}(0)=0$,
$L_{2h+1}'(0)=(2h+1)L_{2h}(0)$ and Stirling's formula.
Write $t=dx$ and
\[
 V_d(t)=\frac{(1+4t/d)^{d/2+1/4}}{2\sqrt\pi\,t^{5/4}} .
\]
The upper energy inequality gives
$|Y(x)|/d^{3/4}\le (1+o(1))V_d(t)$, with the prefactor uniform
in $x>0$. In each bounded lobe choose a point where $u'=0$, possible
by Rolle. At this point the lower bound is
$(1+o(1))V_d(t)\sqrt{q(\pi/2)/q(\theta)}$.
The square root is bounded below by an absolute positive constant
throughout every such lobe. Indeed, $|L_d|\le1$ on $[-1,1]$, so
Markov's polynomial inequality gives $|L_d'|\le d^2$ there, and
$L_d(\cos\theta)>0$ for $0<\theta\le1/d$. The square root converges to
one uniformly when $t$ ranges over a compact subinterval of $(0,\infty)$.

\emph{Counting lobes.}
Write $u=r\sin\psi$, $u'=nr\cos\psi$, with continuous phase and $r>0$.
The differential equation gives
$\psi'=n+\sin^2\psi/(4n\sin^2\theta)$. On any subinterval of
$[1/d,\pi/2]$ the integral of the second term is $O(1)$, by integrating
$\csc^2\theta/(4n)$. Thus the number of all zero gaps in
$[\theta_1,\theta_2]\subseteq[1/d,\pi/2]$ is
$n(\theta_2-\theta_1)/\pi+O(1)$, with half as many positive lobes up
to $O(1)$. The same estimate counts lobes meeting the interval instead
of those contained in it. Each angle gap has length at most $\pi/n$.

On a compact $t$-interval bounded away from zero the gap widths in $t$
therefore tend uniformly to zero, and $V_d(t)\to A_0(t)$ uniformly.
The maximum of each positive lobe, divided by $d^{3/4}$, is consequently
$A_0(t)+o(1)$ uniformly for any $t$ in it. For a sequence of chosen
extrema of $u$ with $t\to0$ or $t\to\infty$, the lower bound for the
normalized lobe maximum diverges. At zero this follows from $t^{-5/4}$.
At infinity, $V_d$ increases past $t=5d/(8d-16)$ and
$A_0(T)\to\infty$ as fixed $T\to\infty$. There are lobes arbitrarily
close to $t=5/8$, the unique minimum of $A_0$. These facts prove (1).

For (2), away from small neighborhoods of $t_-$ and $t_+$, the lobes
with maxima below $M$ are precisely those in $(t_-,t_+)$ for large $d$.
In a fixed $t$-interval $[a,b]\subset(0,\infty)$,
$n(\theta_b-\theta_a)/\pi=(2/\pi)(\sqrt b-\sqrt a)\sqrt d+o(\sqrt d)$.
Halving this count to take positive lobes and doubling for nonreal
roots proves (2), by shrinking the endpoint neighborhoods. For $c=c_{\rm num}$
only an arbitrarily small neighborhood of the minimum can contribute.

For (3) put $k_0=(\alpha-3/4)/2$. For any fixed $0<k<k_0$, throughout
$1\le t\le k\log d$ the upper bound $V_d(t)=O(d^{2k})$ shows that
every positive lobe contained there has maximum below $M$. For any
fixed $k>k_0$, the value
$V_d(k\log d)=d^{2k+o(1)}$ and the monotonicity of $V_d$ show that
every positive lobe whose chosen extremum is above $k\log d$ has
maximum greater than $M$. The angle at $t=k\log d$ is
$(2+o(1))\sqrt{k\log d/d}$, while the angle at $t=1$ is $O(d^{-1/2})$.
Apply the gap count to these upper and lower bounds and let $k\to k_0$.
\end{proof}

Within the family of complete graphs with a matching deleted, $d=169$
is the first dimension in which the numerator can have a nonreal
zero. For $3\le d\le168$ the program \texttt{matching\_threshold.py}
exhibits $\lfloor d/2\rfloor$ disjoint positive intervals with
alternating endpoint signs for $\Gamma_{d,M}(-x)$ at the maximum allowed
matching size. It uses trigonometric estimates only to suggest rational
endpoints; their ordering and the signs are tested by integer arithmetic.
The lobe comparison above covers all smaller matchings when $d\ge4$.
Dimensions $d\le3$ have gamma degree at most one. Across the smaller
dimensions and the check at $(d,M)=(169,84)$ there are 7,139 disjoint
root enclosures, counted within their respective polynomials.

At $d=169$ the last admissible deletion is decisive. A matching of $84$
edges gives $84$ real gamma roots and no nonreal roots. Adding the
$85$th edge to the matching leaves $82$ real roots and produces one
nonreal conjugate pair, hence four nonreal $h^*$-roots. The program
\texttt{numerator\_check.py} uses primitive signed Euclidean
remainders. Its Sturm variations at $-\infty,+\infty$ are $(84,0)$
and $(83,1)$, respectively, and both terminal remainders are nonzero
constants. An even shorter certificate suffices to show failure of
real-rootedness at $M=85$. For $P=\Gamma_{169,85}$ and
$x=-1/218$, exact rational arithmetic gives
\[
 -13348 < P'(x)^2-P(x)P''(x) < -13347.
\]
A real polynomial with only real roots satisfies
$(P')^2-PP''\ge0$ on $\mathbb R$, by differentiating its logarithmic
derivative away from the roots and extending by continuity. The checker evaluates $P,P',P''$ together by Horner's recurrence in
exact rational arithmetic.
